\chapter{Advantages and Limitations}
\label{ch:adv}

\section{Advantages}

\begin{enumerate}
  \item \textbf{Long context at low cost.} The model supports contexts of up
        to one million tokens. Because 40 of its 48 sequence-mixing layers are
        Mamba-2 blocks with a constant-size state, only 8 layers need a KV
        cache, and the cost of generating each new token grows far more slowly
        with context length than in a pure Transformer.
  \item \textbf{High throughput.} It runs up to 2.2$\times$ faster than
        GPT-OSS-120B and up to 7.5$\times$ faster than Qwen3.5-122B at
        comparable accuracy, in a long-generation setting.
  \item \textbf{Large capacity, small active size.} Only about 12 billion of
        the 120 billion parameters are used for any one token, so the compute
        per token is that of a much smaller model.
  \item \textbf{Hardware-aware MoE.} LatentMoE reduces the memory traffic and
        the routing traffic of each expert by a factor of four and uses the
        saving to run more experts, improving accuracy at about the same cost.
  \item \textbf{Built-in speculative decoding.} The shared-weight MTP heads act
        as an internal draft model, with an average of 3.45 tokens accepted
        per verification step and no separate draft model to deploy.
  \item \textbf{Efficient training and deployment precision.} Pretraining in
        4-bit NVFP4 was stable over 25 trillion tokens, and the 4-bit
        deployment checkpoint keeps a median of 99.8\% of full-precision
        accuracy.
  \item \textbf{Openness.} Base, post-trained and quantized checkpoints, much
        of the training data and the training recipe are released openly, so
        the work can be inspected, reproduced in part and built upon.
\end{enumerate}

\section{Limitations}

\begin{enumerate}
  \item \textbf{Memory footprint.} MoE reduces computation, not storage. All
        120 billion parameters must still be held in memory even though only a
        tenth are used per token, so the hardware requirements remain high.
  \item \textbf{Serving complexity.} Expert routing makes deployment more
        complicated than for a dense model: different tokens need different
        experts, and in multi-GPU serving tokens must be routed between
        devices. Combining Mamba, attention, MoE and MTP also means the serving
        software must support all four efficiently.
  \item \textbf{Sensitivity of the recurrent state.} Compressing the Mamba
        state causes rounding errors that build up over long outputs
        (Chapter~\ref{ch:quant}). The stochastic-rounding fix works, but it
        depends on hardware support that is specific to recent NVIDIA GPUs, and
        the underlying effect was observed rather than removed by training.
  \item \textbf{Hardware dependence of the headline speed.} The 2.2$\times$
        figure applies to the NVFP4 checkpoint on Blackwell GPUs. On hardware
        without native FP4 support the advantage is smaller.
  \item \textbf{Accuracy gap to the strongest competitor.} The model is still
        behind Qwen3.5-122B on several reasoning and agentic benchmarks,
        including GPQA, HLE, SWE-Bench and TauBench.
  \item \textbf{Lossy memory in most layers.} A Mamba state is a summary. Exact
        recall over long distances relies on just eight attention layers, and
        the published results do not include an ablation showing how accuracy
        changes with the number or placement of those layers at this scale.
  \item \textbf{Cost of training.} Pretraining on 25 trillion tokens followed
        by RL across 21 environments on thousands of GPUs is far beyond the
        reach of most organisations; the open release lowers the barrier to
        \emph{using} the model, not to reproducing it from scratch.
  \item \textbf{Vendor-reported results.} All figures in this report come from
        the developers' own technical reports. Independent evaluation on other
        hardware and other serving frameworks would be needed to confirm how
        general the advantages are.
\end{enumerate}

\section{Comparison at a Glance}

Table~\ref{tab:compare} places \NS\ beside a pure Transformer MoE model of the
same class.

\begin{table}[htbp]
\centering
\small
\caption{Hybrid Mamba--Transformer MoE compared with a pure Transformer MoE}
\label{tab:compare}
\begin{tabularx}{\linewidth}{@{}P{3.6cm}LL@{}}
\toprule
\textbf{Aspect} & \textbf{Pure Transformer MoE} & \textbf{\NS\ (hybrid)} \\
\midrule
Sequence mixing & Attention in every block & Mamba-2 in 40 blocks, attention in 8 \\
Generation memory & KV cache in every block & KV cache in 8 layers, fixed state elsewhere \\
Feed-forward part & Standard MoE at full width & LatentMoE in a 1024-dimensional latent space \\
Faster decoding & External draft model, if any & Built-in MTP heads \\
Deployment risk & Well understood & Recurrent state needs careful quantization \\
Tooling maturity & Very mature & Newer, fewer optimised runtimes \\
\bottomrule
\end{tabularx}
\end{table}
